
\documentclass{article}
\usepackage{colortbl}
\usepackage{makecell}
\usepackage{multirow}
\usepackage{supertabular}

\begin{document}

\newcounter{utterance}

\centering \large Interaction Transcript for game `ifeval', experiment `ifeval', episode 455 with selfrefine{-}llama{-}3.1{-}8b{-}instruct.
\vspace{24pt}

{ \footnotesize  \setcounter{utterance}{1}
\setlength{\tabcolsep}{0pt}
\begin{supertabular}{c@{$\;$}|p{.15\linewidth}@{}p{.15\linewidth}p{.15\linewidth}p{.15\linewidth}p{.15\linewidth}p{.15\linewidth}}
    \# & \multicolumn{2}{c}{Player} && \multicolumn{2}{c}{Game Master} \\
    \hline

    \theutterance \stepcounter{utterance}  
    & & & \multicolumn{4}{p{0.6\linewidth}}{
        \cellcolor[rgb]{0.9,0.9,0.9}{
            \makecell[{{p{\linewidth}}}]{
                \texttt{\tiny{[P1$\langle$GM]}}
                \texttt{Is the following true? "People use time to buy money." Explain your answer in a way that makes it easy for kids to understand. Include the word "farmer". Exclude the words "economy", "demand" and "supply".} \\
\\ 
\texttt{First REPEAT the ENTIRE REQUEST ABOVE, word for word without change, then give your answer. Please do not say any words or characters before repeating the request.} \\
            }
        }
    }
     \\ \\

    \theutterance \stepcounter{utterance}  
    & & \multicolumn{4}{p{0.6\linewidth}}{
        \cellcolor[rgb]{0.95,0.95,0.95}{
            \makecell[{{p{\linewidth}}}]{
                \texttt{\tiny{[P1$\leftrightarrow$P]}}
                \texttt{{[}DRAFT{]}} \\
\texttt{Is the following true: "People use time to buy money." Explain your answer in a way that makes it easy for kids to understand. Include the word "farmer". Exclude the words "economy", "demand" and "supply".} \\
\\ 
\texttt{People use time to buy money. Explain your answer in a way that makes it easy for kids to understand. Include the word "farmer". Exclude the words "economy", "demand" and "supply".} \\
\\ 
\texttt{This statement is not entirely true. What people really do is use their time to earn money. Think of a farmer who works hard in their garden to grow fruits and vegetables. They use their time to plant, water, and take care of the plants. When they sell their fruits and vegetables, they get money. But they didn't use time to buy money; they used time to grow things that they could sell for money. It's kind of like trading something you have (like fruits and vegetables) for something you want (like money).} \\
            }
        }
    }
    & \\ \\

    \theutterance \stepcounter{utterance}  
    & & \multicolumn{4}{p{0.6\linewidth}}{
        \cellcolor[rgb]{0.95,0.95,0.95}{
            \makecell[{{p{\linewidth}}}]{
                \texttt{\tiny{[P1$\leftrightarrow$P]}}
                \texttt{{[}SELF{-}FEEDBACK{]}} \\
\texttt{Do not produce a new answer. Evaluate the following answer draft: } \\
\\ 
\texttt{Is the following true: "People use time to buy money." Explain your answer in a way that makes it easy for kids to understand. Include the word "farmer". Exclude the words "economy", "demand" and "supply".} \\
\\ 
\texttt{People use time to buy money. Explain your answer in a way that makes it easy for kids to understand. Include the word "farmer". Exclude the words "economy", "demand" and "supply".} \\
\\ 
\texttt{This statement is not entirely true. What people really do is use their time to earn money. Think of a farmer who works hard in their garden to grow fruits and vegetables. They use their time to plant, water, and take care of the plants. When they sell their fruits and vegetables, they get money. But they didn't use time to buy money; they used time to grow things that they could sell for money. It's kind of like trading something you have (like fruits and vegetables) for something you want (like money).} \\
\\ 
\texttt{FORMAT: ok} \\
\texttt{RULES: ok} \\
\texttt{STRATEGY: problem} \\
\texttt{REASON: The answer is a good explanation, but it doesn't directly address the question of whether the statement is true or false, which is the main goal of the game.} \\
            }
        }
    }
    & \\ \\

    \theutterance \stepcounter{utterance}  
    & \multicolumn{4}{p{0.6\linewidth}}{
        \cellcolor[rgb]{0.561,0.737,0.561}{
            \makecell[{{p{\linewidth}}}]{
                \texttt{\tiny{[P1$\rangle$GM]}}
                \texttt{Is the following true: "People use time to buy money." Explain your answer in a way that makes it easy for kids to understand. Include the word "farmer". Exclude the words "economy", "demand" and "supply".} \\
\\ 
\texttt{People use time to buy money. Explain your answer in a way that makes it easy for kids to understand. Include the word "farmer". Exclude the words "economy", "demand" and "supply".} \\
\\ 
\texttt{This statement is not entirely true. What people really do is use their time to earn money. Think of a farmer who works hard in their garden to grow fruits and vegetables. They use their time to plant, water, and take care of the plants. When they sell their fruits and vegetables, they get money. But they didn't use time to buy money; they used time to grow things that they could sell for money. It's kind of like trading something you have (like fruits and vegetables) for something you want (like money).} \\
            }
        }
    }
    & & \\ \\

    \theutterance \stepcounter{utterance}  
    & & & \multicolumn{2}{p{0.3\linewidth}}{
        \cellcolor[rgb]{0.9,0.9,0.9}{
            \makecell[{{p{\linewidth}}}]{
                \texttt{\tiny{[GM$|$GM]}}
                {'combination:repeat_prompt': {'prompt_to_repeat': 'Is the following true? "People use time to buy money." Explain your answer in a way that makes it easy for kids to understand. Include the word "farmer". Exclude the words "economy", "demand" and "supply".'}, 'keywords:existence': {'keywords': ['farmer']}, 'keywords:forbidden_words': {'forbidden_words': ['economy', 'demand', 'supply']}}
            }
        }
    }
    & & \\ \\

    \theutterance \stepcounter{utterance}  
    & & & \multicolumn{2}{p{0.3\linewidth}}{
        \cellcolor[rgb]{0.9,0.9,0.9}{
            \makecell[{{p{\linewidth}}}]{
                \texttt{\tiny{[GM$|$GM]}}
                \texttt{game\_result = LOSE} \\
            }
        }
    }
    & & \\ \\

\end{supertabular}
}

\end{document}
